\documentclass[12pt,a4paper]{article}
\usepackage[T1]{fontenc}
\usepackage[utf8]{inputenc}
\usepackage{lmodern}
\usepackage{amsmath,amssymb,amsthm,mathtools}
\usepackage{geometry}
\geometry{margin=1in}
\usepackage{microtype}
\usepackage{hyperref}
\usepackage{xcolor}
\usepackage{enumitem}
\usepackage{fancyhdr}
\usepackage{bookmark}

\usepackage{centernot}

\DeclareMathOperator{\rank}{rank}

\DeclareMathOperator{\Deck}{Deck}

\usepackage{tikz-cd}

\usepackage{tocloft}
\setlength{\cftsubsecnumwidth}{3.2em}

\newcommand{\Z}{\mathbb{Z}}
\newcommand{\Q}{\mathbb{Q}}
\newcommand{\RP}{\mathbb{RP}}
\newcommand{\Chi}{\chi}
\newcommand{\del}{\partial}

\newtheorem{prop}{Proposition}[section]

\newtheorem{theorem}{Theorem}
\newtheorem{proposition}{Proposition}
\newtheorem{lemma}{Lemma}
\newtheorem{example}{Example}


\hypersetup{
	colorlinks=true,
	linkcolor=blue!50!black,
	urlcolor=blue!50!black,
	citecolor=blue!50!black,
	pdftitle={Theory of Algebraic Topology 6},
	pdfauthor={},
	pdfsubject={Solved problems and theory notes in commutative algebra},
	pdfcreator={OpenAI}
}

\setlist[itemize]{topsep=4pt,itemsep=2pt,parsep=0pt,partopsep=0pt}
\setlength{\parskip}{0.55em}
\setlength{\parindent}{0pt}

\setcounter{secnumdepth}{2}


\setcounter{tocdepth}{2}

\pagestyle{fancy}


\fancyhf{}
\lhead{Theory of Algebraic Topology 6}
\rhead{\thepage}
\renewcommand{\headrulewidth}{0.4pt}

\newtheoremstyle{notespace}{6pt}{6pt}{\normalfont}{} {\bfseries}{.}{0.5em}{}
\theoremstyle{notespace}
\newtheorem{definition}{Definition}[section]
\newtheorem{remark}{Remark}[section]

\DeclareMathOperator{\Spec}{Spec}
\DeclareMathOperator{\Ass}{Ass}
\DeclareMathOperator{\Ann}{Ann}
\DeclareMathOperator{\Supp}{Supp}
\DeclareMathOperator{\Hom}{Hom}
\DeclareMathOperator{\Min}{Min}
\DeclareMathOperator{\Ker}{Ker}
\DeclareMathOperator{\Frac}{Frac}

\DeclareMathOperator{\depth}{depth}
\DeclareMathOperator{\Tor}{Tor}
\DeclareMathOperator{\Ext}{Ext}

\DeclareMathOperator{\coker}{coker}
\DeclareMathOperator{\im}{im}

\newcommand{\longdash}{\mathrel{\text{---}}}

\title{\vspace{-1.5em}\bfseries Theory of Algebraic Topology 6}
\author{}
\date{\today}

\begin{document}
	\maketitle
	\thispagestyle{plain}
		\begin{center}
		\begin{minipage}{0.92\textwidth}
			Lecture notes with worked examples in algebraic topology
		\end{minipage}
	\end{center}
		\pdfbookmark[1]{Contents}{contents}\tableofcontents
	\clearpage

\section{Homology: definitions, intuition, and step-by-step examples}

Homology is an algebraic tool for detecting and measuring the presence of ``holes'' in a topological space. Very roughly,
\[
H_0 \text{ measures connected components,}
\qquad
H_1 \text{ measures loops,}
\qquad
H_2 \text{ measures shells or void boundaries,}
\]
and higher \(H_n\) measure higher-dimensional holes.

The central idea is simple:
\begin{itemize}
	\item take formal sums of simplices or cells,
	\item define their boundaries,
	\item look at those chains whose boundary is zero,
	\item then quotient out by those chains that are themselves boundaries of higher-dimensional chains.
\end{itemize}
Thus homology records closed geometric objects that do \emph{not} come from filling something in.

\subsection{Chain groups and boundary maps}

Let \(X\) be a simplicial complex or a CW complex. For each \(n\geq 0\), the group \(C_n(X)\) of \(n\)-chains is the free abelian group generated by the oriented \(n\)-simplices (or \(n\)-cells) of \(X\).

For example, if \(X\) has edges \(e_1,e_2,e_3\), then
\[
C_1(X)=\mathbb{Z}e_1\oplus \mathbb{Z}e_2\oplus \mathbb{Z}e_3.
\]
If \(X\) has vertices \(v_1,v_2\), then
\[
C_0(X)=\mathbb{Z}v_1\oplus \mathbb{Z}v_2.
\]

There is a boundary homomorphism
\[
\partial_n:C_n(X)\longrightarrow C_{n-1}(X).
\]
For an oriented edge \([v_0,v_1]\),
\[
\partial_1([v_0,v_1])=v_1-v_0.
\]
For an oriented triangle \([v_0,v_1,v_2]\),
\[
\partial_2([v_0,v_1,v_2])
=
[v_1,v_2]-[v_0,v_2]+[v_0,v_1].
\]

A fundamental fact is
\[
\partial_{n-1}\circ \partial_n=0,
\]
that is, the boundary of a boundary is always zero.

\subsection{Cycles, boundaries, and homology groups}

An \(n\)-chain \(c\in C_n(X)\) is called an \(n\)-cycle if
\[
\partial_n(c)=0.
\]
The subgroup of \(n\)-cycles is
\[
Z_n(X)=\ker(\partial_n).
\]

An \(n\)-chain is called an \(n\)-boundary if it is the boundary of some \((n+1)\)-chain:
\[
B_n(X)=\operatorname{im}(\partial_{n+1}).
\]

Since \(\partial_{n-1}\circ \partial_n=0\), every boundary is a cycle:
\[
B_n(X)\subseteq Z_n(X).
\]

The \(n\)-th homology group is defined by
\[
H_n(X)=Z_n(X)/B_n(X)=\ker(\partial_n)/\operatorname{im}(\partial_{n+1}).
\]

So homology keeps track of cycles modulo those cycles that already bound something.

\subsection{Geometric intuition}

The geometric meaning is:

\begin{itemize}
	\item \(H_0(X)\) detects connected components.
	\item \(H_1(X)\) detects essential loops that do not bound a 2-dimensional region.
	\item \(H_2(X)\) detects closed surfaces that do not bound a 3-dimensional region inside the space.
\end{itemize}

For instance:
\begin{itemize}
	\item a loop around a disk is trivial in homology, because it is the boundary of the disk;
	\item a loop around a circle is nontrivial, because there is no 2-dimensional region in the circle filling it in;
	\item the surface of a sphere is nontrivial in \(H_2\), because it is a closed 2-dimensional shell.
\end{itemize}

\subsection{Example: the interval \([0,1]\)}

Take the interval with vertices \(v_0,v_1\) and one edge \(e=[v_0,v_1]\). Then
\[
C_1=\mathbb{Z}e,
\qquad
C_0=\mathbb{Z}v_0\oplus \mathbb{Z}v_1.
\]
The boundary map is
\[
\partial_1(e)=v_1-v_0.
\]
There are no 2-simplices, so \(\partial_2=0\).

To compute \(H_1\), let \(ae\in C_1\). Then
\[
\partial_1(ae)=a(v_1-v_0).
\]
This is zero only if \(a=0\). Hence
\[
Z_1=0,
\qquad
B_1=0,
\qquad
H_1([0,1])=0.
\]

Now
\[
Z_0=C_0=\mathbb{Z}v_0\oplus \mathbb{Z}v_1,
\]
and
\[
B_0=\operatorname{im}(\partial_1)=\langle v_1-v_0\rangle.
\]
Therefore
\[
H_0([0,1])
\cong
\frac{\mathbb{Z}v_0\oplus \mathbb{Z}v_1}{\langle v_1-v_0\rangle}
\cong
\mathbb{Z}.
\]

Thus
\[
H_0([0,1])\cong \mathbb{Z},
\qquad
H_1([0,1])=0.
\]
This reflects the fact that the interval is connected and has no 1-dimensional hole.

\subsection{Example: the circle \(S^1\)}

A convenient CW decomposition of \(S^1\) has one 0-cell \(v\) and one 1-cell \(e\), attached at both ends to \(v\). Then
\[
C_1=\mathbb{Z}e,
\qquad
C_0=\mathbb{Z}v.
\]
The boundary is
\[
\partial_1(e)=v-v=0.
\]
There are no 2-cells.

Hence
\[
Z_1=\ker(\partial_1)=\mathbb{Z}e,
\qquad
B_1=0,
\]
so
\[
H_1(S^1)\cong \mathbb{Z}.
\]
Also,
\[
H_0(S^1)\cong \mathbb{Z}.
\]

Therefore
\[
H_0(S^1)\cong \mathbb{Z},
\qquad
H_1(S^1)\cong \mathbb{Z}.
\]

The generator of \(H_1(S^1)\) is the class of going once around the circle.

\subsection{Example: the sphere \(S^2\)}

A simple CW structure on \(S^2\) has one 0-cell and one 2-cell, with no 1-cells. Thus
\[
C_2\cong \mathbb{Z},
\qquad
C_1=0,
\qquad
C_0\cong \mathbb{Z}.
\]
All boundary maps are zero, so
\[
H_2(S^2)\cong \mathbb{Z},
\qquad
H_1(S^2)=0,
\qquad
H_0(S^2)\cong \mathbb{Z}.
\]

Thus the sphere has one connected component, no essential loops, and one essential 2-dimensional shell.

More generally,
\[
H_k(S^n)=
\begin{cases}
	\mathbb{Z}, & k=0,n,\\
	0, & \text{otherwise}.
\end{cases}
\]

\subsection{Example: the figure-eight \(S^1\vee S^1\)}

Let \(X=S^1\vee S^1\), the wedge of two circles. It has one vertex \(v\) and two 1-cells \(a,b\). Then
\[
C_1=\mathbb{Z}a\oplus \mathbb{Z}b,
\qquad
C_0=\mathbb{Z}v.
\]
Since both loops start and end at the same vertex,
\[
\partial_1(a)=0,
\qquad
\partial_1(b)=0.
\]
There are no 2-cells, so
\[
Z_1=C_1\cong \mathbb{Z}^2,
\qquad
B_1=0,
\]
hence
\[
H_1(S^1\vee S^1)\cong \mathbb{Z}^2.
\]
Also,
\[
H_0(S^1\vee S^1)\cong \mathbb{Z}.
\]

So
\[
H_0(S^1\vee S^1)\cong \mathbb{Z},
\qquad
H_1(S^1\vee S^1)\cong \mathbb{Z}^2.
\]

Geometrically, there are two independent essential loops.

\subsection{Example: the torus \(T^2\)}

The torus has homology groups
\[
H_0(T^2)\cong \mathbb{Z},
\qquad
H_1(T^2)\cong \mathbb{Z}^2,
\qquad
H_2(T^2)\cong \mathbb{Z}.
\]

This can be understood geometrically:
\begin{itemize}
	\item \(H_0\cong \mathbb{Z}\): the torus is connected;
	\item \(H_1\cong \mathbb{Z}^2\): there are two basic independent loops, the meridian and the longitude;
	\item \(H_2\cong \mathbb{Z}\): the whole oriented torus surface gives a fundamental 2-dimensional class.
\end{itemize}

\subsection{Reduced homology}

Reduced homology \(\widetilde{H}_n(X)\) is a small modification of ordinary homology that removes the extra \(\mathbb{Z}\) in degree \(0\) for a connected nonempty space.

For a connected nonempty space \(X\),
\[
\widetilde{H}_0(X)=0,
\qquad
\widetilde{H}_n(X)\cong H_n(X)\quad \text{for } n\geq 1.
\]

For example,
\[
\widetilde{H}_0(S^1)=0,
\qquad
\widetilde{H}_1(S^1)\cong \mathbb{Z},
\]
and
\[
\widetilde{H}_0(S^2)=0,
\qquad
\widetilde{H}_2(S^2)\cong \mathbb{Z}.
\]

Reduced homology is especially useful for wedge sums and suspensions.

\subsection{Free part, torsion, and Betti numbers}

Every finitely generated abelian group can be written in the form
\[
\mathbb{Z}^r\oplus T,
\]
where \(T\) is a finite abelian group. The summand \(\mathbb{Z}^r\) is called the \emph{free part}, and \(T\) is the \emph{torsion part}.

For example,
\[
\mathbb{Z}^2\oplus \mathbb{Z}/3\oplus \mathbb{Z}/12
\]
has free part \(\mathbb{Z}^2\) and torsion part \(\mathbb{Z}/3\oplus \mathbb{Z}/12\).

If
\[
H_n(X)\cong \mathbb{Z}^r\oplus T,
\]
then the integer \(r\) is called the \(n\)-th Betti number:
\[
b_n(X)=\operatorname{rank} H_n(X).
\]

The free part corresponds to genuinely independent holes, while torsion reflects finite-order phenomena.

\subsection{Rational homology}

If we tensor homology with \(\mathbb{Q}\), we get
\[
H_n(X;\mathbb{Q})=H_n(X;\mathbb{Z})\otimes \mathbb{Q}.
\]
This removes torsion, because
\[
(\mathbb{Z}/m)\otimes \mathbb{Q}=0
\]
for every \(m\geq 1\), while
\[
\mathbb{Z}\otimes \mathbb{Q}\cong \mathbb{Q}.
\]

Hence, if
\[
H_n(X;\mathbb{Z})\cong \mathbb{Z}^r\oplus T,
\]
then
\[
H_n(X;\mathbb{Q})\cong \mathbb{Q}^r.
\]

So rational homology sees only the free part of integral homology.

\subsection{Example: the projective plane \(\mathbb{RP}^2\)}

A standard computation gives
\[
H_0(\mathbb{RP}^2)\cong \mathbb{Z},
\qquad
H_1(\mathbb{RP}^2)\cong \mathbb{Z}/2,
\qquad
H_2(\mathbb{RP}^2)=0.
\]
Therefore over the rationals,
\[
H_0(\mathbb{RP}^2;\mathbb{Q})\cong \mathbb{Q},
\qquad
H_1(\mathbb{RP}^2;\mathbb{Q})=0,
\qquad
H_2(\mathbb{RP}^2;\mathbb{Q})=0.
\]

Thus the first homology of \(\mathbb{RP}^2\) is purely torsion, and rational homology does not detect it.

The relation
\[
H_1(\mathbb{RP}^2)\cong \mathbb{Z}/2
\]
means that there is a nontrivial 1-dimensional homology class \(\alpha\) such that
\[
2\alpha=0
\qquad \text{but} \qquad \alpha\neq 0.
\]

\subsection{Suspension}

The suspension \(\Sigma X\) of a space \(X\) is obtained from \(X\times [0,1]\) by collapsing all of \(X\times \{0\}\) to one point and all of \(X\times \{1\}\) to another point. Geometrically, one adds a north pole and a south pole and stretches \(X\) between them.

Examples:
\[
\Sigma S^0\cong S^1,
\qquad
\Sigma S^1\cong S^2,
\qquad
\Sigma S^n\cong S^{n+1}.
\]

A basic theorem says that reduced homology shifts by one under suspension:
\[
\widetilde{H}_n(\Sigma X)\cong \widetilde{H}_{n-1}(X).
\]

For example, since
\[
\widetilde{H}_1(S^1)\cong \mathbb{Z},
\]
it follows that
\[
\widetilde{H}_2(\Sigma S^1)\cong \mathbb{Z}.
\]
Since \(\Sigma S^1\cong S^2\), this agrees with
\[
\widetilde{H}_2(S^2)\cong \mathbb{Z}.
\]

\subsection{Wedge sums}

If \(X\) and \(Y\) are pointed spaces, their wedge sum \(X\vee Y\) is obtained by taking the disjoint union and identifying one base point of \(X\) with one base point of \(Y\).

Examples:
\[
S^1\vee S^1
\]
is the figure-eight, and
\[
S^2\vee S^2
\]
is two spheres touching at one point.

For reduced homology,
\[
\widetilde{H}_n(X\vee Y)\cong \widetilde{H}_n(X)\oplus \widetilde{H}_n(Y).
\]

For example,
\[
\widetilde{H}_2(S^2\vee S^2)\cong \mathbb{Z}\oplus \mathbb{Z}.
\]

\subsection{A quotient example: \(S^2/S^1\)}

If one collapses the equator \(S^1\subset S^2\) to a point, the quotient is homeomorphic to
\[
S^2/S^1 \cong S^2\vee S^2.
\]
Indeed, the upper hemisphere becomes a sphere after collapsing its boundary, and the lower hemisphere does the same; the two meet at the collapse point.

Therefore
\[
\widetilde{H}_2(S^2/S^1)\cong \mathbb{Z}\oplus \mathbb{Z}.
\]

\subsection{Exact sequences in homology}

A sequence
\[
A\xrightarrow{f} B\xrightarrow{g} C
\]
is called exact at \(B\) if
\[
\operatorname{im}(f)=\ker(g).
\]
This means that the elements of \(B\) coming from \(A\) are exactly those killed by the map to \(C\).

One of the most important tools in algebraic topology is the long exact sequence of a pair:
\[
\cdots \to H_n(A)\to H_n(X)\to H_n(X,A)\to H_{n-1}(A)\to H_{n-1}(X)\to \cdots
\]
for a pair \((X,A)\).

Another central tool is the Mayer--Vietoris sequence. If \(X=U\cup V\), then there is a long exact sequence
\[
\cdots \to H_n(U\cap V)\to H_n(U)\oplus H_n(V)\to H_n(X)\to H_{n-1}(U\cap V)\to \cdots
\]
which is a homological analogue of inclusion--exclusion.

\subsection{Homology and the fundamental group}

For a path-connected space \(X\), the first homology group is the abelianization of the fundamental group:
\[
H_1(X;\mathbb{Z})\cong \pi_1(X)^{ab}.
\]

For example,
\[
\pi_1(S^1\vee S^1)\cong F_2,
\]
the free group on two generators, and its abelianization is
\[
F_2^{ab}\cong \mathbb{Z}^2.
\]
Thus
\[
H_1(S^1\vee S^1)\cong \mathbb{Z}^2.
\]

So \(H_1\) remembers loops, but only in an abelianized way; it forgets the noncommutative structure of \(\pi_1\).

\subsection{Table of standard examples}

\[
\begin{array}{c|ccc}
	X & H_0(X) & H_1(X) & H_2(X)\\
	\hline
	[0,1] & \mathbb{Z} & 0 & 0\\
	S^1 & \mathbb{Z} & \mathbb{Z} & 0\\
	S^2 & \mathbb{Z} & 0 & \mathbb{Z}\\
	S^1\vee S^1 & \mathbb{Z} & \mathbb{Z}^2 & 0\\
	T^2 & \mathbb{Z} & \mathbb{Z}^2 & \mathbb{Z}\\
	\mathbb{RP}^2 & \mathbb{Z} & \mathbb{Z}/2 & 0
\end{array}
\]

Over \(\mathbb{Q}\), this becomes
\[
\begin{array}{c|ccc}
	X & H_0(X;\mathbb{Q}) & H_1(X;\mathbb{Q}) & H_2(X;\mathbb{Q})\\
	\hline
	S^1 & \mathbb{Q} & \mathbb{Q} & 0\\
	S^2 & \mathbb{Q} & 0 & \mathbb{Q}\\
	T^2 & \mathbb{Q} & \mathbb{Q}^2 & \mathbb{Q}\\
	\mathbb{RP}^2 & \mathbb{Q} & 0 & 0
\end{array}
\]

\subsection{How to compute homology in practice}

A practical method for computing homology is:

\begin{enumerate}
	\item Choose a decomposition of the space: simplicial, cellular, or a decomposition suited for Mayer--Vietoris.
	\item Write down the chain groups \(C_n\).
	\item Write the boundary maps \(\partial_n\).
	\item Compute the cycles:
	\[
	Z_n=\ker(\partial_n).
	\]
	\item Compute the boundaries:
	\[
	B_n=\operatorname{im}(\partial_{n+1}).
	\]
	\item Form the quotient:
	\[
	H_n=Z_n/B_n.
	\]
	\item Interpret the result geometrically in terms of connected components, loops, shells, and torsion phenomena.
\end{enumerate}

\subsection{Summary}

The philosophy of homology can be compressed into one formula:
\[
H_n(X)=\frac{\text{\(n\)-cycles}}{\text{\(n\)-boundaries}}
=
\frac{\ker(\partial_n)}{\operatorname{im}(\partial_{n+1})}.
\]

Thus homology measures those closed \(n\)-dimensional objects that do not already arise as boundaries of higher-dimensional ones. It is one of the basic algebraic invariants in topology and provides a systematic language for discussing connectedness, loops, surfaces, and higher-dimensional holes.

\section{Problem 6 Torus endomorphism induced by an integer matrix}

\textbf{Problem statement.}
Let \(A\) be a \(2\times 2\) matrix with entries in \(\mathbb Z\). Show that the linear map
\[
A:\mathbb R^2\to \mathbb R^2
\]
preserves the equivalence relation
\[
(a,b)\sim (a',b')
\quad\Longleftrightarrow\quad
(a-a',\,b-b')\in \mathbb Z^2,
\]
and therefore induces a continuous map
\[
f_A:T\to T
\]
on the torus
\[
T=\mathbb R^2/\mathbb Z^2.
\]
Compute the induced homomorphisms
\[
(f_A)_*:H_k(T)\to H_k(T)
\]
on the homology of \(T\).

\bigskip

\textbf{Solution.}
Recall that the torus is the quotient
\[
T=\mathbb R^2/\mathbb Z^2,
\]
where two points of \(\mathbb R^2\) are identified if they differ by an element of the lattice \(\mathbb Z^2\).

We first show that \(A\) respects this equivalence relation. Suppose
\[
(a,b)\sim (a',b').
\]
Then
\[
(a-a',\,b-b')\in \mathbb Z^2.
\]
Writing \(x=(a,b)\) and \(y=(a',b')\), this means \(x-y\in \mathbb Z^2\). Since \(A\) has integer entries, it sends integer vectors to integer vectors:
\[
A(\mathbb Z^2)\subseteq \mathbb Z^2.
\]
Hence
\[
A(x-y)=Ax-Ay\in \mathbb Z^2,
\]
so \(Ax\sim Ay\). Therefore \(A\) preserves the quotient relation and induces a well-defined map
\[
f_A:T\to T,
\qquad
f_A([x])=[Ax].
\]
Because the quotient map \(\pi:\mathbb R^2\to T\) is continuous and \(A\) is continuous, the induced map \(f_A\) is continuous.

We now compute the induced map on homology.

For the torus,
\[
H_0(T)\cong \mathbb Z,\qquad
H_1(T)\cong \mathbb Z^2,\qquad
H_2(T)\cong \mathbb Z,
\]
and \(H_k(T)=0\) for all \(k\ge 3\).

\subsection*{Action on \(H_0(T)\)}
Since \(T\) is connected, any continuous self-map \(T\to T\) induces the identity on \(H_0\). Thus
\[
(f_A)_0=\mathrm{id}_{\mathbb Z}.
\]

\subsection*{Action on \(H_1(T)\)}
Choose the standard generators of \(H_1(T)\cong \mathbb Z^2\) represented by the two coordinate circles:
\[
\alpha(t)=[(t,0)],
\qquad
\beta(t)=[(0,t)],
\qquad t\in \mathbb R/\mathbb Z.
\]
Write
\[
A=
\begin{pmatrix}
	a & b\\
	c & d
\end{pmatrix}.
\]
Then
\[
A(t,0)=(at,ct),
\qquad
A(0,t)=(bt,dt).
\]
Therefore the induced loops satisfy
\[
f_A\circ \alpha \sim a\alpha+c\beta,
\qquad
f_A\circ \beta \sim b\alpha+d\beta.
\]
Hence, with respect to the basis \([\alpha],[\beta]\) of \(H_1(T)\),
\[
(f_A)_1=
\begin{pmatrix}
	a & b\\
	c & d
\end{pmatrix}
=A.
\]

\subsection*{Action on \(H_2(T)\)}
Since \(H_2(T)\cong \mathbb Z\), the induced map on \(H_2\) must be multiplication by the degree of \(f_A\). For a torus map induced by the linear map \(A\), the degree is
\[
\deg(f_A)=\det(A).
\]
Thus
\[
(f_A)_2=\times \det(A).
\]

So the induced maps on homology are:
\[
(f_A)_*:H_0(T)\to H_0(T)
\quad\text{is}\quad
\mathrm{id}_{\mathbb Z},
\]
\[
(f_A)_*:H_1(T)\to H_1(T)
\quad\text{is}\quad
A,
\]
\[
(f_A)_*:H_2(T)\to H_2(T)
\quad\text{is multiplication by}\quad
\det(A),
\]
and for \(k\ge 3\),
\[
H_k(T)=0.
\]

Thus the complete answer is
\[
(f_A)_0=\mathrm{id},\qquad
(f_A)_1=A,\qquad
(f_A)_2=\times \det(A).
\]



\section{Problem 7 Connected sum of an orientable and a nonorientable surface}

\textbf{Problem statement.}
For triangulated surfaces \(X\) and \(Y\), let \(X\#Y\) be the surface obtained by removing a \(2\)-simplex from each and gluing the two boundary circles formed by \(\partial\Delta^2\). Let \(\Sigma_g\) denote the closed orientable surface of genus \(g\), and let \(S_n\) denote the closed nonorientable surface of genus \(n\). Use the Mayer--Vietoris sequence to compute the homology of
\[
\Sigma_g\# S_n,
\]
and deduce that
\[
\Sigma_g\# S_n\cong S_{n+2g}.
\]

\bigskip

\textbf{Solution.}
Let
\[
M=\Sigma_g\# S_n.
\]
By construction, we may write
\[
M=U\cup V,
\]
where
\[
U=\Sigma_g\setminus \operatorname{int}(D^2),
\qquad
V=S_n\setminus \operatorname{int}(D^2),
\]
and the intersection \(U\cap V\) is a collar neighborhood of the glued boundary circle. Therefore
\[
U\cap V\simeq S^1.
\]

We now describe the homology of the pieces.

\subsection*{Step 1: homology of \(U\), \(V\), and \(U\cap V\)}
The space \(U\) is an orientable surface of genus \(g\) with one boundary component. It deformation retracts onto a wedge of \(2g\) circles, so
\[
H_0(U)\cong \mathbb Z,
\qquad
H_1(U)\cong \mathbb Z^{2g},
\qquad
H_i(U)=0\ \text{for }i\ge 2.
\]

Similarly, \(V\) is a nonorientable surface of genus \(n\) with one boundary component. It deformation retracts onto a wedge of \(n\) circles, so
\[
H_0(V)\cong \mathbb Z,
\qquad
H_1(V)\cong \mathbb Z^{n},
\qquad
H_i(V)=0\ \text{for }i\ge 2.
\]

Finally, since \(U\cap V\simeq S^1\),
\[
H_0(U\cap V)\cong \mathbb Z,
\qquad
H_1(U\cap V)\cong \mathbb Z,
\qquad
H_i(U\cap V)=0\ \text{for }i\ge 2.
\]

\subsection*{Step 2: the Mayer--Vietoris sequence}
The relevant part of the Mayer--Vietoris sequence for \(M=U\cup V\) is
\[
0\to H_2(M)\to H_1(U\cap V)
\overset{\alpha}{\longrightarrow}
H_1(U)\oplus H_1(V)
\longrightarrow H_1(M)
\longrightarrow H_0(U\cap V)
\longrightarrow H_0(U)\oplus H_0(V)
\longrightarrow H_0(M)\to 0.
\]
Substituting the homology groups above gives
\[
0\to H_2(M)\to \mathbb Z
\overset{\alpha}{\longrightarrow}
\mathbb Z^{2g}\oplus \mathbb Z^n
\longrightarrow H_1(M)
\longrightarrow \mathbb Z
\longrightarrow \mathbb Z\oplus \mathbb Z
\longrightarrow H_0(M)\to 0.
\]

Since \(U\), \(V\), and \(U\cap V\) are all connected, the map
\[
H_0(U\cap V)\to H_0(U)\oplus H_0(V)
\]
is the diagonal map \(1\mapsto (1,1)\), hence injective. It follows that
\[
H_0(M)\cong \mathbb Z,
\]
and exactness reduces the computation of \(H_1(M)\) and \(H_2(M)\) to the short exact sequence
\[
0\to H_2(M)\to \mathbb Z
\overset{\alpha}{\longrightarrow}
\mathbb Z^{2g}\oplus \mathbb Z^n
\longrightarrow H_1(M)\to 0.
\]

Thus the key point is to determine the map \(\alpha\).

\subsection*{Step 3: identify the boundary class}
The generator of \(H_1(U\cap V)\cong \mathbb Z\) is represented by the common boundary circle.

\paragraph{Its image in \(H_1(U)\).}
For the orientable surface \(U=\Sigma_g\setminus \operatorname{int}(D^2)\), the boundary loop is homologous to zero. One way to see this is that
\[
\pi_1(U)=
\left\langle
a_1,b_1,\dots,a_g,b_g
\right\rangle
\]
with the boundary represented by the word
\[
[a_1,b_1]\cdots [a_g,b_g].
\]
After abelianization, every commutator becomes trivial, so the boundary class maps to \(0\in H_1(U)\).

\paragraph{Its image in \(H_1(V)\).}
For the nonorientable surface \(V=S_n\setminus \operatorname{int}(D^2)\), one has a standard presentation
\[
\pi_1(V)=\langle x_1,\dots,x_n\rangle,
\]
and the boundary loop is represented by
\[
x_1^2x_2^2\cdots x_n^2.
\]
After abelianization, this becomes
\[
2x_1+2x_2+\cdots+2x_n
\in H_1(V)\cong \mathbb Z^n.
\]

Hence the map
\[
\alpha:H_1(U\cap V)\to H_1(U)\oplus H_1(V)
\]
is given by
\[
\alpha(1)=\bigl(0,\;2x_1+2x_2+\cdots+2x_n\bigr).
\]

Now \(x_1+\cdots+x_n\) is a primitive element of \(\mathbb Z^n\), so after a change of basis in \(\mathbb Z^n\), the map \(\alpha\) becomes
\[
\mathbb Z\to \mathbb Z^{2g}\oplus \mathbb Z^n,
\qquad
1\mapsto (0,2,0,\dots,0).
\]
This map is injective. Therefore
\[
H_2(M)=0.
\]

Moreover,
\[
H_1(M)\cong
\bigl(\mathbb Z^{2g}\oplus \mathbb Z^n\bigr)\big/\langle (0,2,0,\dots,0)\rangle.
\]
Thus
\[
H_1(M)\cong \mathbb Z^{2g+n-1}\oplus \mathbb Z/2.
\]

So we have proved
\[
H_0(\Sigma_g\# S_n)\cong \mathbb Z,
\]
\[
H_1(\Sigma_g\# S_n)\cong \mathbb Z^{2g+n-1}\oplus \mathbb Z/2,
\]
\[
H_2(\Sigma_g\# S_n)=0,
\]
and all higher homology groups vanish.

\subsection*{Step 4: identify the surface}
For the closed nonorientable surface \(S_m\) of genus \(m\), the homology groups are
\[
H_0(S_m)\cong \mathbb Z,
\qquad
H_1(S_m)\cong \mathbb Z^{m-1}\oplus \mathbb Z/2,
\qquad
H_2(S_m)=0.
\]
If we put
\[
m=n+2g,
\]
then
\[
H_1(S_{n+2g})\cong \mathbb Z^{n+2g-1}\oplus \mathbb Z/2,
\]
which agrees exactly with the homology computed above for \(\Sigma_g\# S_n\).

Thus \(\Sigma_g\# S_n\) has the same homology as \(S_{n+2g}\). By the classification theorem for closed connected surfaces, the connected sum of an orientable genus-\(g\) surface with a nonorientable genus-\(n\) surface is the nonorientable genus-\((n+2g)\) surface. Therefore
\[
\Sigma_g\# S_n\cong S_{n+2g}.
\]

So the final result is
\[
H_0(\Sigma_g\# S_n)\cong \mathbb Z,
\qquad
H_1(\Sigma_g\# S_n)\cong \mathbb Z^{n+2g-1}\oplus \mathbb Z/2,
\qquad
H_2(\Sigma_g\# S_n)=0,
\]
and hence
\[
\boxed{\Sigma_g\# S_n\cong S_{n+2g}.}
\]

\section{Deep explanation of the torus: quotient, CW structure, fundamental group, and homology}

The torus is one of the central examples in topology and algebraic topology. It can be understood in several equivalent ways:

\begin{itemize}
	\item as the product \(S^1\times S^1\),
	\item as a square with opposite sides identified,
	\item as the quotient \(\mathbb R^2/\mathbb Z^2\),
	\item as a CW complex with one \(0\)-cell, two \(1\)-cells, and one \(2\)-cell,
	\item as the closed orientable surface of genus \(1\).
\end{itemize}

Each of these descriptions reveals a different aspect of the torus.

\subsection{The torus as a product of circles}

A very clean definition of the torus is
\[
T^2=S^1\times S^1.
\]
Thus a point of the torus is an ordered pair of points on two circles. Intuitively, one circle measures motion in one circular direction, and the second circle measures motion in another independent circular direction.

This already suggests that the torus has two basic loops:

\begin{itemize}
	\item one loop going around the ``tube'',
	\item one loop going around the ``hole''.
\end{itemize}

These two loops will later generate both the fundamental group and the first homology group.

\subsection{The square model of the torus}

A standard geometric model for the torus is
\[
T^2=[0,1]\times[0,1]/\sim,
\]
where the equivalence relation identifies opposite sides:
\[
(x,0)\sim (x,1),
\qquad
(0,y)\sim (1,y).
\]

So we begin with a square and impose the rule that:

\begin{itemize}
	\item the bottom edge is glued to the top edge,
	\item the left edge is glued to the right edge.
\end{itemize}

This gives the torus.

It is helpful to imagine the gluing in two stages.

\paragraph{Stage 1.}
Glue the left edge to the right edge. The square becomes a cylinder.

\paragraph{Stage 2.}
Now glue the two circular boundary components of the cylinder together. The result is a torus.

So the torus is the quotient space obtained by identifying opposite sides of a square.

\subsection{The torus as \(\mathbb R^2/\mathbb Z^2\)}

Another fundamental model is
\[
T^2=\mathbb R^2/\mathbb Z^2.
\]
This means that two points \((x,y)\) and \((x',y')\) of \(\mathbb R^2\) are identified if and only if
\[
(x-x',\,y-y')\in \mathbb Z^2.
\]
Equivalently,
\[
(x,y)\sim (x+m,y+n)
\qquad\text{for all }m,n\in\mathbb Z.
\]

In other words, translating by one unit horizontally or vertically does not change the corresponding point of the torus.

\paragraph{Geometric picture.}
The plane \(\mathbb R^2\) is tiled by unit squares. Each unit square is another copy of the same fundamental domain. When one passes from one square to a neighboring square by an integer translation, one returns to the same point of the torus.

Thus the torus may be viewed as one unit square together with the rule that its opposite sides are identified.

\paragraph{Why this is important.}
This description shows immediately that \(\mathbb R^2\) is the universal cover of the torus, and it makes the action of integer matrices on the torus especially transparent.

\subsection{The torus as a CW complex}

The torus admits a very economical CW structure:
\[
T^2=e^0\cup e^1_a\cup e^1_b\cup e^2.
\]
That is:

\begin{itemize}
	\item one \(0\)-cell,
	\item two \(1\)-cells,
	\item one \(2\)-cell.
\end{itemize}

\paragraph{The \(0\)-skeleton.}
Start with one point \(v\).

\paragraph{The \(1\)-skeleton.}
Attach two intervals by gluing both endpoints of each interval to the point \(v\). This produces two loops, traditionally called \(a\) and \(b\). The resulting \(1\)-skeleton is
\[
T^{(1)}=S^1\vee S^1,
\]
the wedge sum of two circles.

\paragraph{The \(2\)-cell.}
Now attach one \(2\)-disk \(D^2\) along its boundary by the attaching map represented by the word
\[
aba^{-1}b^{-1}.
\]
This means that the boundary circle of the disk travels:

\begin{itemize}
	\item once along \(a\),
	\item then once along \(b\),
	\item then back along \(a\),
	\item then back along \(b\).
\end{itemize}

This attaching word is exactly the topological encoding of the square model of the torus.

\subsection{Why the boundary word is \(aba^{-1}b^{-1}\)}

Take the square model of the torus and walk once around the boundary of the square counterclockwise.

Label:

\begin{itemize}
	\item the bottom edge by \(a\),
	\item the right edge by \(b\).
\end{itemize}

Then:

\begin{itemize}
	\item the top edge is identified with the bottom edge, but is traversed in the opposite direction, so it contributes \(a^{-1}\),
	\item the left edge is identified with the right edge, but is traversed in the opposite direction, so it contributes \(b^{-1}\).
\end{itemize}

Therefore the full boundary word is
\[
aba^{-1}b^{-1}.
\]

This word gives the attaching map of the \(2\)-cell.

\subsection{The fundamental group of the torus}

From the CW structure, one obtains the standard presentation
\[
\pi_1(T^2)=\langle a,b\mid aba^{-1}b^{-1}=1\rangle.
\]
The relation
\[
aba^{-1}b^{-1}=1
\]
is equivalent to
\[
ab=ba.
\]
Thus the two generators commute, and so
\[
\pi_1(T^2)\cong \mathbb Z^2.
\]

This agrees with the product description:
\[
T^2=S^1\times S^1.
\]
Since
\[
\pi_1(S^1)\cong \mathbb Z,
\]
we get
\[
\pi_1(T^2)\cong \pi_1(S^1)\times \pi_1(S^1)\cong \mathbb Z\times \mathbb Z=\mathbb Z^2.
\]

So a loop on the torus is classified, up to homotopy, by two integers: its winding numbers in the two basic directions.

\subsection{The two basic loops}

Let us write the torus as
\[
T^2=S^1\times S^1.
\]
Then two standard loops are
\[
a(t)=(e^{2\pi i t},1),
\qquad
b(t)=(1,e^{2\pi i t}),
\qquad t\in [0,1].
\]
These generate the fundamental group and also the first homology group.

Geometrically:

\begin{itemize}
	\item \(a\) runs around one circular direction,
	\item \(b\) runs around the other circular direction.
\end{itemize}

Any class in \(\pi_1(T^2)\) or \(H_1(T^2)\) is determined by an integer pair
\[
(m,n)\in \mathbb Z^2,
\]
meaning that the loop winds \(m\) times in one direction and \(n\) times in the other.

\subsection{Homology of the torus}

The homology groups of the torus are
\[
H_0(T^2)\cong \mathbb Z,
\qquad
H_1(T^2)\cong \mathbb Z^2,
\qquad
H_2(T^2)\cong \mathbb Z,
\]
and
\[
H_k(T^2)=0
\qquad\text{for all }k\ge 3.
\]

\paragraph{Interpretation of \(H_0\).}
Since the torus is connected,
\[
H_0(T^2)\cong \mathbb Z.
\]

\paragraph{Interpretation of \(H_1\).}
The first homology group records the independent \(1\)-dimensional cycles. The torus has two essential independent loops, represented by \(a\) and \(b\). Hence
\[
H_1(T^2)\cong \mathbb Z^2.
\]

\paragraph{Interpretation of \(H_2\).}
The torus is a closed orientable surface, so it has a fundamental class in dimension \(2\). Therefore
\[
H_2(T^2)\cong \mathbb Z.
\]

\subsection{Cellular homology of the torus}

Using the CW structure, the cellular chain groups are
\[
C_2(T^2)\cong \mathbb Z,
\qquad
C_1(T^2)\cong \mathbb Z^2,
\qquad
C_0(T^2)\cong \mathbb Z.
\]

So the cellular chain complex is
\[
0\longrightarrow \mathbb Z \xrightarrow{\partial_2} \mathbb Z^2 \xrightarrow{\partial_1} \mathbb Z \longrightarrow 0.
\]

We now compute the boundary maps.

\paragraph{The map \(\partial_1\).}
Each \(1\)-cell is a loop based at the same vertex, so its initial and terminal points coincide. Hence
\[
\partial_1=0.
\]

\paragraph{The map \(\partial_2\).}
The \(2\)-cell is attached by the word
\[
aba^{-1}b^{-1}.
\]
In cellular homology, this contributes
\[
a+b-a-b=0.
\]
Thus
\[
\partial_2=0.
\]

Therefore the chain complex becomes
\[
0\longrightarrow \mathbb Z \xrightarrow{0} \mathbb Z^2 \xrightarrow{0} \mathbb Z \longrightarrow 0.
\]

Hence
\[
H_2(T^2)=\ker(\partial_2)\cong \mathbb Z,
\]
\[
H_1(T^2)=\ker(\partial_1)/\operatorname{im}(\partial_2)\cong \mathbb Z^2,
\]
\[
H_0(T^2)=\mathbb Z.
\]

So the cellular computation matches the geometric intuition exactly.

\subsection{Orientability of the torus}

The torus is orientable. Intuitively, one can move continuously around the torus while keeping a consistent notion of orientation.

This is in contrast with nonorientable surfaces such as:

\begin{itemize}
	\item the Möbius band,
	\item the Klein bottle,
	\item \(\mathbb{RP}^2\).
\end{itemize}

Because the torus is closed and orientable, it has nontrivial top homology:
\[
H_2(T^2)\cong \mathbb Z.
\]
For a closed nonorientable surface, the top homology over \(\mathbb Z\) is \(0\).

\subsection{Euler characteristic of the torus}

From the CW structure we immediately obtain
\[
\chi(T^2)=\#\{\text{\(0\)-cells}\}-\#\{\text{\(1\)-cells}\}+\#\{\text{\(2\)-cells}\}.
\]
Since the torus has:

\begin{itemize}
	\item one \(0\)-cell,
	\item two \(1\)-cells,
	\item one \(2\)-cell,
\end{itemize}
we get
\[
\chi(T^2)=1-2+1=0.
\]

This agrees with the general formula for the closed orientable surface \(\Sigma_g\) of genus \(g\):
\[
\chi(\Sigma_g)=2-2g.
\]
For \(g=1\), this gives
\[
\chi(\Sigma_1)=2-2=0.
\]
Thus the torus is precisely the orientable surface of genus \(1\):
\[
T^2=\Sigma_1.
\]

\subsection{Curves on the torus and slopes}

In the model
\[
T^2=\mathbb R^2/\mathbb Z^2,
\]
closed curves on the torus can be understood by looking at lines in the plane.

A line of slope \((m,n)\) in \(\mathbb R^2\) projects to a closed curve on the torus when the slope is rational. In that case the corresponding homology class is
\[
(m,n)\in H_1(T^2)\cong \mathbb Z^2.
\]

If \(\gcd(m,n)=1\), the curve is primitive and does not wrap multiple times around a simpler curve. If \(\gcd(m,n)>1\), then the curve traces a simpler primitive curve several times.

If the slope is irrational, the projected image does not close; instead, it winds densely through the torus.

This is one of the most beautiful features of the quotient model \(\mathbb R^2/\mathbb Z^2\).

\subsection{The universal cover of the torus}

The universal covering map is
\[
p:\mathbb R^2\to T^2=\mathbb R^2/\mathbb Z^2.
\]
Every point of the torus has infinitely many lifts in \(\mathbb R^2\), differing by integer lattice translations.

The deck transformation group is
\[
\mathbb Z^2,
\]
acting by
\[
(m,n)\cdot (x,y)=(x+m,y+n).
\]

This explains geometrically why
\[
\pi_1(T^2)\cong \mathbb Z^2.
\]
The fundamental group is identified with the group of deck transformations of the universal cover.

\subsection{Maps from the torus to itself induced by integer matrices}

One especially elegant property of the torus is that any integer matrix
\[
A\in M_2(\mathbb Z)
\]
induces a continuous self-map
\[
f_A:T^2\to T^2.
\]
Indeed, if \(x-y\in \mathbb Z^2\), then
\[
A(x-y)\in \mathbb Z^2
\]
because \(A\) has integer entries. So the equivalence relation defining \(\mathbb R^2/\mathbb Z^2\) is preserved.

Thus
\[
f_A([x])=[Ax]
\]
is well defined.

Its effect on homology is especially simple:

\[
(f_A)_0=\mathrm{id}
\quad\text{on }H_0(T^2)\cong \mathbb Z,
\]
\[
(f_A)_1=A
\quad\text{on }H_1(T^2)\cong \mathbb Z^2,
\]
\[
(f_A)_2=\times \det(A)
\quad\text{on }H_2(T^2)\cong \mathbb Z.
\]

So the matrix itself acts on the first homology, while its determinant gives the degree on the top homology.

\subsection{Comparison with other surfaces}

It is useful to compare the torus with the sphere and the real projective plane.

\paragraph{Sphere \(S^2\).}
\[
\pi_1(S^2)=0,
\qquad
H_1(S^2)=0,
\qquad
H_2(S^2)\cong \mathbb Z.
\]

\paragraph{Torus \(T^2\).}
\[
\pi_1(T^2)\cong \mathbb Z^2,
\qquad
H_1(T^2)\cong \mathbb Z^2,
\qquad
H_2(T^2)\cong \mathbb Z.
\]

\paragraph{Projective plane \(\mathbb{RP}^2\).}
\[
\pi_1(\mathbb{RP}^2)\cong \mathbb Z/2,
\qquad
H_1(\mathbb{RP}^2)\cong \mathbb Z/2,
\qquad
H_2(\mathbb{RP}^2)=0
\quad\text{over }\mathbb Z.
\]

Thus the torus differs from \(\mathbb{RP}^2\) in two essential ways:

\begin{itemize}
	\item the torus is orientable,
	\item the torus has free first homology \(\mathbb Z^2\), whereas \(\mathbb{RP}^2\) has torsion \(\mathbb Z/2\).
\end{itemize}

\subsection{Summary}

The torus is the same topological space viewed in many equivalent ways:
\[
T^2=S^1\times S^1
=[0,1]^2/\sim
=\mathbb R^2/\mathbb Z^2
=\Sigma_1.
\]

Its basic algebraic-topological invariants are
\[
\pi_1(T^2)\cong \mathbb Z^2,
\qquad
H_0(T^2)\cong \mathbb Z,
\qquad
H_1(T^2)\cong \mathbb Z^2,
\qquad
H_2(T^2)\cong \mathbb Z,
\qquad
\chi(T^2)=0.
\]

The square model explains its quotient structure, the CW model explains its cell attachment structure, the universal cover \(\mathbb R^2\) explains its fundamental group, and the product model \(S^1\times S^1\) explains its two independent circular directions.

\bigskip

\textbf{Solution.}
The torus is a closed orientable surface of genus \(1\), and it may be described equivalently as \(S^1\times S^1\), as a square with opposite sides identified, or as \(\mathbb R^2/\mathbb Z^2\). Its CW structure has one \(0\)-cell, two \(1\)-cells, and one \(2\)-cell attached by the word \(aba^{-1}b^{-1}\). From this one obtains
\[
\pi_1(T^2)\cong \mathbb Z^2,
\qquad
H_0(T^2)\cong \mathbb Z,
\qquad
H_1(T^2)\cong \mathbb Z^2,
\qquad
H_2(T^2)\cong \mathbb Z,
\]
and
\[
\chi(T^2)=0.
\]
The universal cover is \(\mathbb R^2\), with deck transformation group \(\mathbb Z^2\), and any integer matrix \(A\in M_2(\mathbb Z)\) induces a self-map of the torus whose action on \(H_1\) is given by \(A\) and whose action on \(H_2\) is multiplication by \(\det(A)\).

\section{Problem 8 Finite-sheeted coverings, triangulations, Euler characteristic, and coverings of surfaces}

\textbf{Statement.}
Let
\[
p:\widetilde X \to X
\]
be a finite-sheeted covering space, and let
\[
h:|K|\to X
\]
be a triangulation. Show that there is an \(r\ge 1\) and a triangulation
\[
g:|L|\to \widetilde X
\]
such that the composition
\[
h^{-1}\circ p\circ g:|L|\to |K^{(r)}|
\]
is a simplicial map.

If \(p\) has \(n\) sheets, show that
\[
\chi(\widetilde X)=n\,\chi(X).
\]
Hence show that \(\Sigma_g\) is a covering space of \(\Sigma_h\) if and only if
\[
\frac{1-g}{1-h}
\]
is an integer.

\bigskip

\textbf{Solution.}
We proceed in three steps.

\medskip

\textbf{Step 1: after subdivision, the covering becomes simplicial.}

Because \(p:\widetilde X\to X\) is a covering map, each point of \(X\) has an evenly covered neighborhood. Since \(h:|K|\to X\) is a triangulation and \(|K|\) is compact on each simplex, after sufficiently many barycentric subdivisions there exists \(r\ge 1\) such that for every simplex \(\sigma\) of \(K^{(r)}\), the set \(h(|\sigma|)\) lies inside some evenly covered open set \(U_\sigma\subset X\).

For each such simplex \(\sigma\), the preimage \(p^{-1}(U_\sigma)\) is a disjoint union of open sets
\[
p^{-1}(U_\sigma)=\bigsqcup_\alpha V_{\sigma,\alpha},
\]
and each restriction
\[
p|_{V_{\sigma,\alpha}}:V_{\sigma,\alpha}\to U_\sigma
\]
is a homeomorphism.

Since \(|\sigma|\) is contractible, each lift of \(h|_{|\sigma|}:|\sigma|\to U_\sigma\) to \(\widetilde X\) lands in one sheet \(V_{\sigma,\alpha}\), and under the inverse of
\[
p|_{V_{\sigma,\alpha}},
\]
it identifies with a copy of \(|\sigma|\). Doing this for every simplex of \(K^{(r)}\), and gluing along common faces, produces a simplicial complex \(L\) together with a triangulation
\[
g:|L|\to \widetilde X
\]
such that
\[
h^{-1}\circ p\circ g:|L|\to |K^{(r)}|
\]
is simplicial.

Thus the covering map may be made simplicial after sufficiently fine subdivision of the triangulation downstairs.

\medskip

\textbf{Step 2: Euler characteristic multiplies by the number of sheets.}

Assume now that \(p\) has exactly \(n\) sheets. Since every simplex of \(|K^{(r)}|\) lies in an evenly covered set, each simplex of \(K^{(r)}\) has exactly \(n\) distinct lifts to simplices of \(L\). Therefore, for each dimension \(q\),
\[
f_q(L)=n\,f_q\bigl(K^{(r)}\bigr),
\]
where \(f_q(\cdot)\) denotes the number of \(q\)-simplices.

Hence
\[
\chi(\widetilde X)
=\sum_{q\ge 0}(-1)^q f_q(L)
=\sum_{q\ge 0}(-1)^q n\,f_q\bigl(K^{(r)}\bigr)
=n\sum_{q\ge 0}(-1)^q f_q\bigl(K^{(r)}\bigr).
\]
Since barycentric subdivision does not change Euler characteristic,
\[
\chi\bigl(|K^{(r)}|\bigr)=\chi(|K|)=\chi(X).
\]
Therefore
\[
\boxed{\chi(\widetilde X)=n\,\chi(X).}
\]

\medskip

\textbf{Step 3: application to orientable closed surfaces.}

Recall that for the closed orientable surface \(\Sigma_m\) of genus \(m\),
\[
\chi(\Sigma_m)=2-2m.
\]
If \(\Sigma_g\to \Sigma_h\) is an \(n\)-sheeted covering, then by the formula just proved,
\[
2-2g=n(2-2h).
\]
Dividing by \(2\),
\[
1-g=n(1-h),
\]
so
\[
\boxed{\frac{1-g}{1-h}=n\in \mathbb Z_{>0}.}
\]
Thus integrality of \((1-g)/(1-h)\) is necessary.

Conversely, assume
\[
1-g=n(1-h)
\]
for some positive integer \(n\), equivalently
\[
g=1+n(h-1).
\]
We explain a standard construction of an \(n\)-sheeted covering
\[
\Sigma_g\to \Sigma_h.
\]

Start with a torus \(T^2\), viewed as \(\mathbb R^2/\mathbb Z^2\). Let \(\mathbb Z/n\) act freely on \(T^2\) by translation in one direction:
\[
[j]\cdot [x,y]=\left[x+\frac{j}{n},\,y\right].
\]
This action is free, and the quotient is again a torus.

Choose \(n\) disjoint open disks in \(T^2\) forming one orbit under this action. Remove them. The resulting surface is a torus with \(n\) boundary components. Now attach to each boundary circle a copy of the surface \(\Sigma_{h-1,1}\), meaning the orientable surface of genus \(h-1\) with one boundary component. Let \(\mathbb Z/n\) cyclically permute these \(n\) attached copies.

The resulting closed orientable surface has genus
\[
1+n(h-1)=g,
\]
and the quotient by the \(\mathbb Z/n\)-action has genus
\[
1+(h-1)=h.
\]
Hence the quotient map gives an \(n\)-sheeted covering
\[
\Sigma_g\to \Sigma_h.
\]

So for \(h\neq 1\),
\[
\boxed{\Sigma_g \text{ covers } \Sigma_h \iff \frac{1-g}{1-h}\in \mathbb Z.}
\]

\medskip

\textbf{Remark.}
When \(h=1\), the denominator \(1-h\) is \(0\), so the displayed ratio is not defined. In that case one uses Euler characteristic directly:
\[
\chi(\Sigma_g)=n\chi(\Sigma_1)=n\cdot 0=0,
\]
hence \(g=1\). Thus only the torus covers the torus.


\section{Problem 9 Coverings of spaces whose universal cover is an even-dimensional sphere}

\textbf{Statement.}
Let
\[
p:S^{2k}\to X
\]
be a covering map, let
\[
G=\pi_1(X,x_0),
\]
and recall that \(G\) then acts freely on \(S^{2k}\). Show that for any
\[
g\in G\setminus\{e\},
\]
the map
\[
g_*:H_{2k}(S^{2k})\to H_{2k}(S^{2k})
\]
is multiplication by \(-1\). Deduce that \(G\) is either trivial or \(\mathbb Z/2\), and that \(\mathbb{RP}^{2k}\) is not a proper covering space of any other space.

\bigskip

\textbf{Solution.}
Since \(S^{2k}\) is simply connected for \(k\ge 1\), the covering
\[
p:S^{2k}\to X
\]
is the universal covering of \(X\). Therefore the fundamental group
\[
G=\pi_1(X,x_0)
\]
identifies with the group of deck transformations of \(p\), acting freely on \(S^{2k}\).

Let \(g\in G\setminus\{e\}\). Then \(g:S^{2k}\to S^{2k}\) is a fixed-point-free homeomorphism.

\medskip

\textbf{Step 1: compute the induced map on top homology.}

For any continuous map
\[
f:S^{2k}\to S^{2k},
\]
the Lefschetz number is
\[
L(f)=\sum_i (-1)^i \operatorname{tr}\bigl(f_*:H_i(S^{2k})\to H_i(S^{2k})\bigr).
\]
The only nonzero homology groups of \(S^{2k}\) are
\[
H_0(S^{2k})\cong \mathbb Z,
\qquad
H_{2k}(S^{2k})\cong \mathbb Z.
\]
Hence
\[
L(f)=\operatorname{tr}(f_*|_{H_0})+\operatorname{tr}(f_*|_{H_{2k}}).
\]
Since \(S^{2k}\) is connected, every self-map induces the identity on \(H_0\), so
\[
\operatorname{tr}(f_*|_{H_0})=1.
\]
On top homology, the induced map is multiplication by the degree:
\[
f_*|_{H_{2k}}=\deg(f).
\]
Therefore
\[
L(f)=1+\deg(f).
\]

Apply this to the deck transformation \(g\). Since \(g\) has no fixed point, the Lefschetz fixed point theorem implies
\[
L(g)=0.
\]
Thus
\[
0=1+\deg(g),
\]
so
\[
\deg(g)=-1.
\]
Equivalently,
\[
\boxed{g_*:H_{2k}(S^{2k})\to H_{2k}(S^{2k}) \text{ is multiplication by } -1.}
\]

\medskip

\textbf{Step 2: deduce that \(G\) is trivial or \(\mathbb Z/2\).}

We now have a homomorphism
\[
\rho:G\to \operatorname{Aut}\bigl(H_{2k}(S^{2k})\bigr)\cong \{\pm 1\},
\]
and every nontrivial element of \(G\) is sent to \(-1\).

Suppose \(g,h\in G\setminus\{e\}\). Then
\[
\rho(gh)=\rho(g)\rho(h)=(-1)(-1)=1.
\]
But the only element of \(G\) that can act by \(+1\) is the identity, because every nontrivial element acts by \(-1\). Hence
\[
gh=e.
\]
In particular, taking \(h=g\), we get
\[
g^2=e.
\]
So every nontrivial element has order \(2\), and there can be at most one nontrivial element. Therefore
\[
\boxed{G\cong \{e\}\ \text{or}\ \mathbb Z/2.}
\]

\medskip

\textbf{Step 3: \(\mathbb{RP}^{2k}\) is not a proper covering of another space.}

Assume for contradiction that there exists a proper covering
\[
q:\mathbb{RP}^{2k}\to Y.
\]
Compose it with the universal double covering
\[
\pi:S^{2k}\to \mathbb{RP}^{2k}.
\]
Then
\[
q\circ \pi:S^{2k}\to Y
\]
is a covering map. Since \(S^{2k}\) is simply connected, it is the universal covering of \(Y\), and thus
\[
\pi_1(Y)
\]
must be either trivial or \(\mathbb Z/2\), by the result of Step 2.

On the other hand,
\[
\pi_1(\mathbb{RP}^{2k})\cong \mathbb Z/2.
\]
For connected coverings, intermediate coverings correspond to subgroups of the fundamental group of the base. Since
\[
q:\mathbb{RP}^{2k}\to Y
\]
is assumed to be a proper covering, the subgroup corresponding to \(\mathbb{RP}^{2k}\) must be a proper subgroup of \(\pi_1(Y)\). But this subgroup is isomorphic to \(\mathbb Z/2\), and neither of the groups
\[
\{e\},\qquad \mathbb Z/2
\]
has a proper subgroup isomorphic to \(\mathbb Z/2\).

This is impossible. Therefore no such proper covering exists. Hence
\[
\boxed{\mathbb{RP}^{2k}\text{ is not a proper covering space of any other space.}}
\]


\section{Problem 10 A simplicial automorphism, its fixed set, and the Lefschetz number}

\textbf{Statement.}
If
\[
f:K\to K
\]
is a simplicial isomorphism, let
\[
X\subset |K|
\]
be the fixed set of \(|f|\), that is,
\[
X=\{x\in |K| \mid |f|(x)=x\}.
\]
Show that the Lefschetz number \(L(f)\) is equal to \(\chi(X)\).

\[
\textit{Hint: barycentrically subdivide \(K\) so that \(X\) is the polyhedron of a sub-simplicial complex.}
\]

\bigskip

\textbf{Solution.}
Let \(K'=\operatorname{Sd}(K)\) be the barycentric subdivision of \(K\). The simplicial isomorphism
\[
f:K\to K
\]
induces a simplicial isomorphism
\[
f':K'\to K'.
\]
Since barycentric subdivision does not change the underlying polyhedron and does not change the induced map on homology up to canonical identification, we have
\[
L(f')=L(f).
\]
Thus it suffices to prove that
\[
L(f')=\chi(X).
\]

\medskip

\textbf{Step 1: describe the fixed set as a subcomplex of the barycentric subdivision.}

Recall that the vertices of \(K'\) are the barycenters \(b_\sigma\) of simplices \(\sigma\) of \(K\). The induced simplicial map \(f'\) acts on vertices by
\[
f'(b_\sigma)=b_{f(\sigma)}.
\]
Therefore a vertex \(b_\sigma\) of \(K'\) is fixed by \(f'\) if and only if
\[
f(\sigma)=\sigma.
\]

A simplex of \(K'\) is a chain
\[
\sigma_0<\sigma_1<\cdots<\sigma_q
\]
of simplices of \(K\). This simplex is fixed pointwise by \(f'\) precisely when every \(\sigma_i\) is fixed setwise by \(f\), because then each vertex \(b_{\sigma_i}\) is fixed.

Hence the collection of all simplices of \(K'\) whose vertices are fixed by \(f'\) forms a subcomplex \(L\subset K'\). Its polyhedron is exactly the fixed-point set:
\[
|L|=X.
\]
Indeed, if a point of \(|K|\) is fixed by \(|f|\), then in the barycentric subdivision it lies in a simplex whose vertices correspond to \(f\)-invariant simplices, and conversely every point of such a simplex is fixed.

Thus
\[
X=|L|
\]
for a subcomplex \(L\subset K'\).

\medskip

\textbf{Step 2: compute the Lefschetz number on the chain level.}

Let \(C_q(K')\) be the simplicial chain group in dimension \(q\), with basis the oriented \(q\)-simplices of \(K'\). The map \(f'\) induces a chain map
\[
f'_q:C_q(K')\to C_q(K').
\]

The trace of \(f'_q\) is determined as follows.

If an oriented \(q\)-simplex \(\tau\) is not carried to itself by \(f'\), then \(\tau\) is sent to a different basis element, so its contribution to the trace is \(0\).

If \(\tau\) is fixed by \(f'\), then every vertex of \(\tau\) is fixed. Hence \(f'\) fixes \(\tau\) pointwise, not merely setwise, and therefore preserves its orientation. So the contribution of such a simplex to the trace is \(+1\).

Thus
\[
\operatorname{tr}(f'_q)=\#\{\text{\(q\)-simplices of }L\}.
\]

Now the Lefschetz number is
\[
L(f')=\sum_{q\ge 0}(-1)^q \operatorname{tr}(f'_q)
=\sum_{q\ge 0}(-1)^q \#\{\text{\(q\)-simplices of }L\}.
\]
But this alternating sum is exactly the Euler characteristic of \(L\):
\[
L(f')=\chi(L).
\]
Since \(X=|L|\), Euler characteristic is preserved under passing from a finite simplicial complex to its polyhedron, so
\[
\chi(L)=\chi(X).
\]
Therefore
\[
L(f)=L(f')=\chi(X).
\]

Hence
\[
\boxed{L(f)=\chi(X).}
\]



	
	\section{Definitions, theory, and examples related to Problems 7--10}
	
	These notes collect the main ideas behind Problems 7--10:
	\begin{itemize}[leftmargin=2em]
		\item connected sums of surfaces and the Mayer--Vietoris sequence,
		\item covering spaces and Euler characteristic,
		\item free actions on even-dimensional spheres and the case of $\RP^{2k}$,
		\item the Lefschetz number and fixed sets of simplicial automorphisms.
	\end{itemize}
	
	Throughout:
	\begin{itemize}[leftmargin=2em]
		\item $\Sigma_g$ denotes the closed orientable surface of genus $g$,
		\item $S_n$ denotes the closed nonorientable surface of genus $n$, that is, the connected sum of $n$ copies of $\RP^2$.
	\end{itemize}
	
	Thus:
	\[
	\Sigma_0=S^2,\qquad \Sigma_1=T^2,\qquad S_1=\RP^2,\qquad S_2=\text{Klein bottle}.
	\]
	
	\subsection*{Basic homology of closed surfaces}
	
	The homology groups of the standard closed surfaces are:
	\[
	H_i(\Sigma_g;\Z)\cong
	\begin{cases}
		\Z & i=0,\\
		\Z^{2g} & i=1,\\
		\Z & i=2,\\
		0 & \text{otherwise},
	\end{cases}
	\]
	and
	\[
	H_i(S_n;\Z)\cong
	\begin{cases}
		\Z & i=0,\\
		\Z^{n-1}\oplus \Z/2 & i=1,\\
		0 & i=2,\\
		0 & \text{otherwise}.
	\end{cases}
	\]
	
	Also:
	\[
	\Chi(\Sigma_g)=2-2g,\qquad \Chi(S_n)=2-n.
	\]
	
	\section{Connected sums of surfaces and Mayer--Vietoris}
	
	\subsection{Connected sum}
	
	\begin{definition}
		Let $X$ and $Y$ be closed surfaces. Their \emph{connected sum} $X\#Y$ is obtained by removing an open disk from each surface and then gluing the two boundary circles together by a homeomorphism.
	\end{definition}
	
	In a triangulated setting, one often removes the interior of a $2$-simplex from each triangulated surface and glues the two copies of the boundary $\partial\Delta^2$.
	
	Geometrically:
	\begin{itemize}[leftmargin=2em]
		\item $X\#S^2\cong X$,
		\item $\Sigma_g\#\Sigma_h\cong \Sigma_{g+h}$,
		\item $S_m\#S_n\cong S_{m+n}$,
		\item $\Sigma_g\#S_n\cong S_{n+2g}$.
	\end{itemize}
	
	The last identity is exactly the phenomenon appearing in Problem 7.
	
	\subsection{Why removing a $2$-simplex is the same as removing a disk}
	
	A $2$-simplex is topologically a closed disk. Its interior is an open disk, and its boundary is a circle:
	\[
	\Delta^2\cong D^2,\qquad \partial \Delta^2\cong S^1.
	\]
	So cutting out a $2$-simplex from a triangulated surface is the combinatorial version of cutting out a disk.
	
	\subsection{Mayer--Vietoris sequence}
	
	\begin{theorem}[Mayer--Vietoris]
		If a space $X$ is written as $X=A\cup B$, where $A$ and $B$ are subspaces with suitable niceness assumptions, then there is a long exact sequence
		\[
		\cdots \to H_n(A\cap B)\to H_n(A)\oplus H_n(B)\to H_n(X)\to H_{n-1}(A\cap B)\to \cdots
		\]
		relating the homology of the pieces and the whole space.
	\end{theorem}
	
	The philosophy is simple: if a space is built from two overlapping pieces, then its homology is determined by the homology of the pieces and their intersection.
	
	\subsection{How Mayer--Vietoris is used for connected sums}
	
	Let
	\[
	X\#Y = A\cup B,
	\]
	where:
	\begin{itemize}[leftmargin=2em]
		\item $A$ is $X$ with an open $2$-simplex removed,
		\item $B$ is $Y$ with an open $2$-simplex removed,
		\item $A\cap B \cong S^1\times (0,1)$, so homotopy equivalent to $S^1$.
	\end{itemize}
	
	Since removing an open disk from a closed surface does not change $H_1$, one has
	\[
	H_1(A)\cong H_1(X),\qquad H_1(B)\cong H_1(Y).
	\]
	Also
	\[
	H_1(A\cap B)\cong H_1(S^1)\cong \Z,\qquad H_2(A)=H_2(B)=0,
	\]
	because $A$ and $B$ now have boundary.
	
	The Mayer--Vietoris sequence in degree $1$ and $2$ becomes
	\[
	0\to H_2(X\#Y)\to H_1(S^1)\to H_1(X)\oplus H_1(Y)\to H_1(X\#Y)\to H_0(S^1)\to \cdots
	\]
	and from this one extracts the homology of the connected sum.
	
	\subsection{Key surface consequence}
	
	\begin{proposition}
		For closed surfaces,
		\[
		\Sigma_g\#S_n\cong S_{n+2g}.
		\]
	\end{proposition}
	
	\textbf{Idea.}
	A handle can be replaced by two crosscaps in the classification of compact surfaces. On the homology level:
	\[
	H_1(\Sigma_g)\cong \Z^{2g},\qquad H_1(S_n)\cong \Z^{n-1}\oplus \Z/2,
	\]
	and the connected sum has first homology
	\[
	H_1(\Sigma_g\#S_n)\cong \Z^{n+2g-1}\oplus \Z/2,
	\]
	which matches the first homology of $S_{n+2g}$.
	
	\subsection{Concrete examples}
	
	\begin{example}[Torus connected sum projective plane]
		\[
		\Sigma_1\#S_1 = T^2\#\RP^2 \cong S_3.
		\]
		So a torus plus one crosscap is the nonorientable surface of genus $3$.
	\end{example}
	
	\begin{example}[Torus connected sum Klein bottle]
		\[
		\Sigma_1\#S_2 \cong S_4.
		\]
		Since $S_2$ is the Klein bottle, this says:
		\[
		T^2\#(\text{Klein bottle})\cong S_4.
		\]
	\end{example}
	
	\begin{example}[Euler characteristic check]
		Since
		\[
		\Chi(X\#Y)=\Chi(X)+\Chi(Y)-2,
		\]
		we get
		\[
		\Chi(\Sigma_g\#S_n)=(2-2g)+(2-n)-2=2-(n+2g)=\Chi(S_{n+2g}).
		\]
		This matches the claimed homeomorphism type.
	\end{example}
	
	\section{Covering spaces, lifted triangulations, and Euler characteristic}
	
	\subsection{Covering maps}
	
	\begin{definition}
		A continuous surjective map
		\[
		p:\widetilde X\to X
		\]
		is a \emph{covering map} if for every $x\in X$ there exists an open neighborhood $U$ of $x$ such that
		\[
		p^{-1}(U)=\bigsqcup_{\alpha} U_\alpha
		\]
		is a disjoint union of open sets, each mapped homeomorphically onto $U$ by $p$.
	\end{definition}
	
	Each set $U$ as above is called \emph{evenly covered}.
	
	\begin{definition}
		If every point of $X$ has exactly $n$ preimages under $p$, then $p$ is called an \emph{$n$-sheeted covering}.
	\end{definition}
	
	\subsection{Deck transformations}
	
	\begin{definition}
		A \emph{deck transformation} of a covering $p:\widetilde X\to X$ is a homeomorphism
		\[
		h:\widetilde X\to \widetilde X
		\]
		such that
		\[
		p\circ h=p.
		\]
	\end{definition}
	
	When $\widetilde X$ is universal, deck transformations correspond to $\pi_1(X)$.
	
	\subsection{Triangulations and lifted simplicial structure}
	
	Suppose $p:\widetilde X\to X$ is a finite-sheeted covering and $h:|K|\to X$ is a triangulation. After sufficiently subdividing $K$, the triangulation lifts to a triangulation
	\[
	g:|L|\to \widetilde X
	\]
	such that the composition
	\[
	h^{-1}\circ p\circ g:|L|\to |K^{(r)}|
	\]
	is simplicial.
	
	This means that upstairs every simplex looks like a disjoint copy of a simplex downstairs. Therefore the number of simplices in each dimension multiplies by the number of sheets.
	
	\subsection{Euler characteristic}
	
	\begin{definition}
		If a finite CW complex or triangulated space has $c_i$ cells (or simplices) in dimension $i$, then its \emph{Euler characteristic} is
		\[
		\Chi(X)=\sum_i (-1)^i c_i.
		\]
	\end{definition}
	
	This is independent of the chosen triangulation or CW decomposition.
	
	\subsection{Multiplicativity under finite covers}
	
	\begin{proposition}
		If $p:\widetilde X\to X$ is an $n$-sheeted covering of finite triangulated spaces, then
		\[
		\Chi(\widetilde X)=n\,\Chi(X).
		\]
	\end{proposition}
	
	\textbf{Reason.}
	If $X$ has $c_i$ simplices in dimension $i$, then $\widetilde X$ has exactly $nc_i$ simplices in dimension $i$, because each simplex lifts to $n$ disjoint copies. Hence
	\[
	\Chi(\widetilde X)=\sum_i (-1)^i (nc_i)=n\sum_i (-1)^i c_i=n\,\Chi(X).
	\]
	
	\subsection{Application to orientable surfaces}
	
	If
	\[
	\Sigma_g\to \Sigma_h
	\]
	is an $n$-sheeted covering, then
	\[
	\Chi(\Sigma_g)=n\,\Chi(\Sigma_h).
	\]
	Since
	\[
	\Chi(\Sigma_g)=2-2g,\qquad \Chi(\Sigma_h)=2-2h,
	\]
	we obtain
	\[
	2-2g=n(2-2h).
	\]
	Dividing by $2$:
	\[
	1-g=n(1-h).
	\]
	Thus
	\[
	n=\frac{1-g}{1-h}.
	\]
	
	So a necessary condition for $\Sigma_g$ to cover $\Sigma_h$ is that
	\[
	\frac{1-g}{1-h}\in \Z.
	\]
	
	Equivalently:
	\[
	g=1+n(h-1).
	\]
	
	\subsection{Interpretation of the formula}
	
	This means that an orientable surface of genus $g$ is an $n$-sheeted cover of an orientable surface of genus $h$ exactly when
	\[
	g=1+n(h-1)
	\]
	for some integer $n\ge 1$.
	
	\begin{remark}
		When $h=1$, the target is the torus. Since $\Chi(\Sigma_1)=0$, any finite-sheeted orientable cover of the torus must also have Euler characteristic $0$, hence must again be a torus.
	\end{remark}
	
	\subsection{Concrete examples}
	
	\begin{example}[Double cover of a genus $2$ surface]
		Let $h=2$ and $n=2$. Then
		\[
		g=1+2(2-1)=3.
		\]
		So $\Sigma_3$ is a $2$-sheeted cover of $\Sigma_2$.
	\end{example}
	
	\begin{example}[Triple cover of a genus $3$ surface]
		Let $h=3$ and $n=3$. Then
		\[
		g=1+3(3-1)=7.
		\]
		So $\Sigma_7$ is a $3$-sheeted cover of $\Sigma_3$.
	\end{example}
	
	\begin{example}[Torus covering torus]
		The map
		\[
		T^2=S^1\times S^1\to S^1\times S^1,\qquad (z,w)\mapsto (z^m,w^n)
		\]
		is an $mn$-sheeted covering.
	\end{example}
	
	\subsection{Group-action viewpoint}
	
	Another way to think about finite-sheeted coverings is through free group actions. If a finite group $G$ acts freely and properly discontinuously on a surface $\Sigma_g$, then the quotient
	\[
	\Sigma_g/G
	\]
	is a surface, and the quotient map
	\[
	\Sigma_g\to \Sigma_g/G
	\]
	is a $|G|$-sheeted covering. Then
	\[
	\Chi(\Sigma_g)=|G|\,\Chi(\Sigma_g/G).
	\]
	
	This is exactly the viewpoint suggested in the hint of Problem 8.
	
	\section{Even-dimensional spheres, degree, and the case of \texorpdfstring{$\RP^{2k}$}{RP2k}}
	
	\subsection{Universal coverings and fundamental groups}
	
	Let
	\[
	p:S^{2k}\to X
	\]
	be a covering map. Then the group
	\[
	G=\pi_1(X,x_0)
	\]
	acts freely on $S^{2k}$ by deck transformations.
	
	Thus every element $g\in G$ gives a homeomorphism
	\[
	g:S^{2k}\to S^{2k}.
	\]
	
	\subsection{Degree of a map on a sphere}
	
	\begin{definition}
		If $f:S^n\to S^n$ is continuous, the induced map
		\[
		f_*:H_n(S^n;\Z)\to H_n(S^n;\Z)\cong \Z
		\]
		must be multiplication by some integer. This integer is called the \emph{degree} of $f$ and is denoted $\deg(f)$.
	\end{definition}
	
	So:
	\[
	f_*(\alpha)=\deg(f)\,\alpha
	\qquad \text{for }\alpha\in H_n(S^n)\cong \Z.
	\]
	
	If $f$ is a homeomorphism, then $\deg(f)=\pm 1$.
	
	\subsection{Why nontrivial deck transformations on $S^{2k}$ have degree $-1$}
	
	The key tool is the Lefschetz number. For a map $f:S^{2k}\to S^{2k}$,
	\[
	H_i(S^{2k};\Q)=
	\begin{cases}
		\Q & i=0,\\
		\Q & i=2k,\\
		0 & \text{otherwise}.
	\end{cases}
	\]
	Therefore
	\[
	L(f)=\operatorname{tr}(f_*|H_0)-\operatorname{tr}(f_*|H_1)+\cdots+\operatorname{tr}(f_*|H_{2k})
	=1+\deg(f).
	\]
	
	Now let $g\neq e$ be a nontrivial deck transformation. Since the action is free, $g$ has no fixed points. By the Lefschetz fixed point theorem, a map with nonzero Lefschetz number must have a fixed point. Hence
	\[
	L(g)=0.
	\]
	So
	\[
	1+\deg(g)=0,
	\]
	and therefore
	\[
	\deg(g)=-1.
	\]
	
	Thus every nontrivial deck transformation acts on
	\[
	H_{2k}(S^{2k})\cong \Z
	\]
	by multiplication by $-1$.
	
	\subsection{Consequence for the deck group}
	
	Let $a,b\in G$ be nontrivial. Then
	\[
	\deg(a)=-1,\qquad \deg(b)=-1.
	\]
	Hence
	\[
	\deg(ab)=\deg(a)\deg(b)=(-1)(-1)=1.
	\]
	But if $ab$ were also nontrivial, then by the same argument it would have degree $-1$, a contradiction. Therefore
	\[
	ab=e.
	\]
	So there can be at most one nontrivial element in $G$.
	
	Hence:
	\[
	G\cong \{e\}\quad\text{or}\quad G\cong \Z/2.
	\]
	
	\begin{proposition}
		If $S^{2k}\to X$ is a covering map, then
		\[
		\pi_1(X)\cong \{e\}\quad\text{or}\quad \Z/2.
		\]
	\end{proposition}
	
	\subsection{The example of projective space}
	
	Real projective space $\RP^n$ is defined as the quotient
	\[
	\RP^n=S^n/\{x\sim -x\}.
	\]
	The quotient map
	\[
	p:S^n\to \RP^n
	\]
	is a $2$-sheeted covering.
	
	For $n=2k$, the nontrivial deck transformation is the antipodal map
	\[
	A(x)=-x.
	\]
	On $S^{2k}$ this has degree
	\[
	\deg(A)=(-1)^{2k+1}=-1.
	\]
	
	So the covering
	\[
	S^{2k}\to \RP^{2k}
	\]
	fits exactly into the previous theory.
	
	\subsection{Why \texorpdfstring{$\RP^{2k}$}{RP2k} is not a proper covering space of another space}
	
	Suppose $\RP^{2k}$ were a proper covering space of some space $Y$. Then we would have coverings
	\[
	S^{2k}\to \RP^{2k}\to Y.
	\]
	Their composition would be a covering
	\[
	S^{2k}\to Y
	\]
	whose deck group would have more than two elements. But we just proved that any deck group acting freely on $S^{2k}$ is either trivial or $\Z/2$. This contradiction shows that $\RP^{2k}$ cannot be a proper covering space of any other space.
	
	\subsection{Concrete examples}
	
	\begin{example}[The sphere $S^2$]
		The antipodal quotient
		\[
		S^2\to \RP^2
		\]
		is a $2$-sheeted cover. The nontrivial deck transformation has degree $-1$.
	\end{example}
	
	\begin{example}[Why this is special to even spheres]
		For odd-dimensional spheres there exist many free cyclic actions. For example,
		\[
		S^1\to S^1,\qquad z\mapsto z^n
		\]
		is an $n$-sheeted cover for every $n\ge 1$.
		So the rigidity in Problem 9 is a special feature of even-dimensional spheres.
	\end{example}
	
	\section{Lefschetz number and fixed-point sets of simplicial isomorphisms}
	
	\subsection{The Lefschetz number}
	
	\begin{definition}
		Let $f:X\to X$ be a continuous map on a space whose rational homology groups are finite-dimensional. The \emph{Lefschetz number} of $f$ is
		\[
		L(f)=\sum_{i\ge 0}(-1)^i\,\operatorname{tr}\!\bigl(f_*:H_i(X;\Q)\to H_i(X;\Q)\bigr).
		\]
	\end{definition}
	
	This is an alternating sum of traces of the induced maps on homology.
	
	\subsection{Lefschetz fixed point theorem}
	
	\begin{theorem}[Lefschetz fixed point theorem]
		If $L(f)\neq 0$, then $f$ has a fixed point.
	\end{theorem}
	
	So the Lefschetz number turns algebraic information on homology into geometric information about fixed points.
	
	\subsection{Simplicial automorphisms}
	
	Let $K$ be a finite simplicial complex, and let
	\[
	f:K\to K
	\]
	be a simplicial isomorphism. It induces a homeomorphism
	\[
	|f|:|K|\to |K|
	\]
	on the polyhedron.
	
	Define the fixed-point set
	\[
	X=\{x\in |K|:\ |f|(x)=x\}.
	\]
	
	Problem 10 states that
	\[
	L(f)=\Chi(X).
	\]
	
	\subsection{Why barycentric subdivision appears}
	
	A fixed set of a simplicial automorphism need not be a subcomplex of the original complex. However, after barycentric subdivision, one can arrange that the fixed set becomes the polyhedron of a genuine simplicial subcomplex.
	
	This is the point of the hint: subdivide $K$ so that the fixed simplices become combinatorially visible.
	
	\subsection{Chain-level intuition}
	
	The simplicial chain groups are
	\[
	C_i(K;\Q),
	\]
	with basis given by oriented $i$-simplices. Since $f$ permutes simplices, it induces linear maps
	\[
	f_\#:C_i(K;\Q)\to C_i(K;\Q).
	\]
	The trace of $f_\#$ in dimension $i$ counts, with signs, the $i$-simplices fixed by $f$.
	
	The Hopf trace formula says:
	\[
	L(f)=\sum_i (-1)^i\operatorname{tr}(f_\#|C_i(K;\Q)).
	\]
	So for a simplicial automorphism, the Lefschetz number can be computed directly from fixed simplices.
	
	If the fixed simplices form a subcomplex with polyhedron $X$, then this alternating sum is exactly the Euler characteristic of $X$:
	\[
	L(f)=\Chi(X).
	\]
	
	\subsection{Main statement}
	
	\begin{proposition}
		Let $f:K\to K$ be a simplicial isomorphism, and let
		\[
		X=\{x\in |K|:\ |f|(x)=x\}
		\]
		be the fixed-point set of $|f|$. Then
		\[
		L(f)=\Chi(X).
		\]
	\end{proposition}
	
	\subsection{Concrete examples}
	
	\begin{example}[Identity map]
		If $f=\mathrm{id}$, then the fixed set is all of $|K|$. Therefore
		\[
		L(f)=\Chi(|K|).
		\]
		This is the simplest case.
	\end{example}
	
	\begin{example}[Reflection of an interval]
		Let $K$ be a subdivided interval and let $f$ be reflection across the midpoint. The fixed-point set is one point. Hence
		\[
		L(f)=\Chi(\{\text{point}\})=1.
		\]
	\end{example}
	
	\begin{example}[Rotation of a triangle]
		Let $K$ be the boundary of a triangle, so $|K|\cong S^1$, and let $f$ rotate the triangle by $120^\circ$. Then there are no fixed points, so
		\[
		X=\varnothing,\qquad \Chi(X)=0.
		\]
		Hence
		\[
		L(f)=0.
		\]
	\end{example}
	
	\begin{example}[Reflection of a square]
		Triangulate a square and reflect it across a diagonal. The fixed-point set is a line segment. Since a segment is contractible,
		\[
		\Chi(X)=1.
		\]
		Therefore
		\[
		L(f)=1.
		\]
	\end{example}
	
	\begin{example}[Antipodal map on $S^{2k}$]
		Let $f(x)=-x$ on $S^{2k}$. There are no fixed points, so $L(f)=0$. On the other hand,
		\[
		L(f)=1+\deg(f)=1+(-1)=0.
		\]
		This matches the general theory perfectly.
	\end{example}
	
	\section{How the four problems fit together}
	
	Problems 7--10 form a coherent block of ideas.
	
	\subsection*{Problem 7}
	Connected sums of surfaces are analyzed by cutting each surface along a disk and gluing along a circle. The Mayer--Vietoris sequence then computes the homology of the result. This shows that
	\[
	\Sigma_g\#S_n
	\]
	has the same homology type as
	\[
	S_{n+2g},
	\]
	and classification of surfaces gives the homeomorphism
	\[
	\Sigma_g\#S_n\cong S_{n+2g}.
	\]
	
	\subsection*{Problem 8}
	A finite covering of triangulated spaces lifts a triangulation, and each simplex downstairs has exactly $n$ lifts upstairs. Therefore Euler characteristic multiplies by the number of sheets:
	\[
	\Chi(\widetilde X)=n\,\Chi(X).
	\]
	Applied to surfaces, this gives the genus formula
	\[
	g=1+n(h-1).
	\]
	
	\subsection*{Problem 9}
	If a group acts freely on an even-dimensional sphere, then every nontrivial element has Lefschetz number $0$, hence degree $-1$. This forces the group to be either trivial or $\Z/2$. As a result,
	\[
	\RP^{2k}
	\]
	cannot be a proper covering space of any other space.
	
	\subsection*{Problem 10}
	For a simplicial isomorphism, the Lefschetz number can be read directly from fixed simplices. After barycentric subdivision, the fixed set becomes a subcomplex, and its Euler characteristic equals the Lefschetz number:
	\[
	L(f)=\Chi(\mathrm{Fix}(|f|)).
	\]
	
	\section{Summary of key formulas}
	
	\[
	\Chi(\Sigma_g)=2-2g,\qquad \Chi(S_n)=2-n.
	\]
	
	\[
	\Chi(\widetilde X)=n\,\Chi(X)\qquad \text{for an $n$-sheeted finite covering.}
	\]
	
	\[
	2-2g=n(2-2h)\qquad \Longleftrightarrow \qquad g=1+n(h-1).
	\]
	
	\[
	L(f)=\sum_i (-1)^i\operatorname{tr}(f_*:H_i(X;\Q)\to H_i(X;\Q)).
	\]
	
	\[
	L(f)=1+\deg(f)\qquad \text{for }f:S^{2k}\to S^{2k}.
	\]
	
	\[
	L(f)=\Chi(\mathrm{Fix}(|f|))\qquad \text{for a simplicial isomorphism }f:K\to K.
	\]
	
	\[
	\Sigma_g\#S_n\cong S_{n+2g}.
	\]
	
\section{Concrete examples for Problems 7--10}

This section collects explicit low-dimensional examples illustrating the ideas behind connected sums, coverings of surfaces, free actions on even-dimensional spheres, and Lefschetz numbers of simplicial maps.

\subsection{Concrete examples for Problem 7: connected sums of surfaces}

Recall the general identity
\[
\Sigma_g \# S_n \cong S_{n+2g},
\]
where \(\Sigma_g\) is the closed orientable surface of genus \(g\), and \(S_n\) is the closed nonorientable surface of genus \(n\).

\subsubsection*{Example 1: \(T^2 \# \RP^2 \cong S_3\)}

Since \(T^2=\Sigma_1\) and \(\RP^2=S_1\), the formula gives
\[
\Sigma_1 \# S_1 \cong S_{1+2\cdot 1}=S_3.
\]
So
\[
T^2 \# \RP^2 \cong S_3.
\]

This means that attaching one torus handle to one projective-plane crosscap produces the nonorientable surface of genus \(3\), that is, the connected sum of three copies of \(\RP^2\).

The Euler characteristic confirms this:
\[
\chi(T^2)=0,\qquad \chi(\RP^2)=1,
\]
hence
\[
\chi(T^2\#\RP^2)=\chi(T^2)+\chi(\RP^2)-2=0+1-2=-1.
\]
On the other hand,
\[
\chi(S_3)=2-3=-1.
\]

The homology also matches:
\[
H_1(S_3;\Z)\cong \Z^2\oplus \Z/2.
\]

\subsubsection*{Example 2: \(T^2 \# K \cong S_4\)}

The Klein bottle is \(S_2\), so
\[
T^2 \# K = \Sigma_1 \# S_2 \cong S_{2+2}=S_4.
\]

Thus
\[
T^2 \# K \cong S_4.
\]

Again the Euler characteristic agrees:
\[
\chi(T^2)=0,\qquad \chi(K)=0,
\]
so
\[
\chi(T^2\#K)=0+0-2=-2.
\]
But
\[
\chi(S_4)=2-4=-2.
\]

This is a good example because it shows clearly that one orientable handle contributes the same topological complexity as two crosscaps.

\subsubsection*{Example 3: \(\Sigma_2 \# \RP^2 \cong S_5\)}

Now take a genus-\(2\) orientable surface and form the connected sum with one projective plane:
\[
\Sigma_2 \# S_1 \cong S_{1+2\cdot 2}=S_5.
\]
So
\[
\Sigma_2 \# \RP^2 \cong S_5.
\]

Here the genus-\(2\) surface contributes two handles, and each handle behaves like two crosscaps in the nonorientable classification. Thus:
\[
2\ \text{handles} \rightsquigarrow 4\ \text{crosscaps},
\]
and together with the original crosscap from \(\RP^2\), one obtains \(5\) crosscaps.

\subsubsection*{Example 4: Euler characteristic formula for connected sums}

For closed surfaces \(X\) and \(Y\),
\[
\chi(X\#Y)=\chi(X)+\chi(Y)-2.
\]

For example, if \(X=\Sigma_g\) and \(Y=S_n\), then
\[
\chi(\Sigma_g\#S_n)=(2-2g)+(2-n)-2=2-(n+2g),
\]
which is exactly the Euler characteristic of \(S_{n+2g}\).

This is a quick numerical check supporting the classification result
\[
\Sigma_g\#S_n\cong S_{n+2g}.
\]

\subsection{Concrete examples for Problem 8: coverings and Euler characteristic}

Recall that if
\[
p:\widetilde X\to X
\]
is an \(n\)-sheeted covering of finite triangulated spaces, then
\[
\chi(\widetilde X)=n\,\chi(X).
\]

For closed orientable surfaces this gives
\[
2-2g=n(2-2h),
\]
so
\[
g=1+n(h-1).
\]

\subsubsection*{Example 1: a covering \(T^2\to T^2\)}

Consider
\[
p:S^1\times S^1\to S^1\times S^1,\qquad p(z,w)=(z^m,w^n),
\]
where \(z,w\in S^1\subset \mathbb C\).

This is an \(mn\)-sheeted covering of the torus by itself.

Indeed, each point \((u,v)\in S^1\times S^1\) has exactly \(m\) choices for an \(m\)-th root of \(u\), and exactly \(n\) choices for an \(n\)-th root of \(v\), so altogether \(mn\) preimages.

The Euler characteristic identity is satisfied automatically:
\[
\chi(T^2)=0,\qquad mn\cdot \chi(T^2)=mn\cdot 0=0.
\]

This is the most explicit standard example of a finite-sheeted cover of a surface.

\subsubsection*{Example 2: a double cover \(\Sigma_3\to \Sigma_2\)}

Take \(h=2\) and \(n=2\). Then
\[
g=1+2(2-1)=3.
\]
So \(\Sigma_3\) is a \(2\)-sheeted covering of \(\Sigma_2\).

Checking with Euler characteristic:
\[
\chi(\Sigma_3)=2-2\cdot 3=-4,
\]
and
\[
2\,\chi(\Sigma_2)=2(2-2\cdot 2)=2(-2)=-4.
\]

Thus the formula works exactly.

\subsubsection*{Example 3: a triple cover \(\Sigma_7\to \Sigma_3\)}

Take \(h=3\) and \(n=3\). Then
\[
g=1+3(3-1)=1+6=7.
\]
So \(\Sigma_7\) can be a \(3\)-sheeted covering of \(\Sigma_3\).

Again:
\[
\chi(\Sigma_7)=2-14=-12,
\]
while
\[
3\,\chi(\Sigma_3)=3(2-6)=3(-4)=-12.
\]

\subsubsection*{Example 4: why \(S^2\) cannot cover a torus or a higher-genus orientable surface}

Suppose there were an \(n\)-sheeted covering
\[
S^2\to \Sigma_h.
\]
Then
\[
\chi(S^2)=n\,\chi(\Sigma_h).
\]
That is,
\[
2=n(2-2h).
\]

If \(h=1\), then \(2-2h=0\), so the right-hand side is \(0\), impossible.

If \(h\ge 2\), then \(2-2h<0\), so the right-hand side is negative, again impossible.

Therefore \(S^2\) cannot be a nontrivial finite cover of any orientable surface of genus \(h\ge 1\).

\subsubsection*{Example 5: the formula \(g=1+n(h-1)\) in practice}

This formula is best remembered as follows:
\[
\boxed{g=1+n(h-1).}
\]

A few sample values are:
\[
\begin{array}{c|c|c}
	h & n & g \\ \hline
	2 & 2 & 3 \\
	2 & 3 & 4 \\
	3 & 2 & 5 \\
	3 & 3 & 7
\end{array}
\]

So:
\[
\Sigma_3\to\Sigma_2 \quad \text{can be 2-sheeted,}
\]
\[
\Sigma_4\to\Sigma_2 \quad \text{can be 3-sheeted,}
\]
\[
\Sigma_5\to\Sigma_3 \quad \text{can be 2-sheeted,}
\]
\[
\Sigma_7\to\Sigma_3 \quad \text{can be 3-sheeted.}
\]

\subsection{Concrete examples for Problem 9: even-dimensional spheres and \(\RP^{2k}\)}

Recall the statement: if
\[
p:S^{2k}\to X
\]
is a covering map, then every nontrivial deck transformation acts on
\[
H_{2k}(S^{2k})\cong \Z
\]
by multiplication by \(-1\). Hence the deck group is either trivial or \(\Z/2\).

\subsubsection*{Example 1: the covering \(S^2\to \RP^2\)}

The real projective plane is
\[
\RP^2=S^2/\{x\sim -x\}.
\]
The quotient map
\[
p:S^2\to \RP^2
\]
is a \(2\)-sheeted covering.

The nontrivial deck transformation is the antipodal map
\[
A:S^2\to S^2,\qquad A(x)=-x.
\]

Its degree is
\[
\deg(A)=(-1)^{2+1}=-1.
\]
Therefore on top homology
\[
A_*:H_2(S^2)\to H_2(S^2)
\]
is multiplication by \(-1\).

This is the model example behind Problem 9.

\subsubsection*{Example 2: the covering \(S^4\to \RP^4\)}

Exactly the same construction works in higher even dimension:
\[
\RP^4=S^4/\{x\sim -x\}.
\]
Again the quotient map
\[
S^4\to \RP^4
\]
is a \(2\)-sheeted covering, with deck group \(\Z/2\).

The nontrivial deck transformation is again antipodal:
\[
x\mapsto -x,
\]
and its degree on \(S^4\) is
\[
(-1)^{4+1}=-1.
\]

So the phenomenon is not special to \(S^2\); it holds for every even sphere.

\subsubsection*{Example 3: why \(S^1\) behaves differently}

On the circle, one has the covering
\[
p:S^1\to S^1,\qquad p(z)=z^n
\]
for every integer \(n\ge 1\).

Its deck group is cyclic of order \(n\), not merely \(\Z/2\).

Thus \(S^1\) admits many nontrivial finite covers, unlike \(S^{2k}\), where the deck group can only be trivial or \(\Z/2\).

This shows that Problem 9 is a special feature of even-dimensional spheres.

\subsubsection*{Example 4: why \(\RP^2\) cannot properly cover another space}

Suppose \(\RP^2\) were a proper covering space of some space \(Y\). Then we would have coverings
\[
S^2\to \RP^2\to Y.
\]
By composing them, we would get a covering
\[
S^2\to Y
\]
with more than \(2\) sheets.

Equivalently, the deck group acting freely on \(S^2\) would have more than two elements.

But Problem 9 shows that a free deck group on an even sphere can only be trivial or \(\Z/2\). This contradiction proves that \(\RP^2\) is not a proper covering space of any other space.

\subsubsection*{Example 5: Lefschetz-number check on the antipodal map}

For the antipodal map \(A:S^{2k}\to S^{2k}\),
\[
L(A)=1+\deg(A)=1+(-1)=0.
\]
This agrees with the fact that \(A\) has no fixed points.

So the algebraic computation of the Lefschetz number exactly matches the geometry.

\subsection{Concrete examples for Problem 10: Lefschetz number and fixed-point sets}

Recall that for a simplicial automorphism \(f:K\to K\), if
\[
X=\{x\in |K|: |f|(x)=x\},
\]
then
\[
L(f)=\chi(X).
\]

\subsubsection*{Example 1: the identity map}

Let \(f=\mathrm{id}_K\). Then every point of \(|K|\) is fixed, so
\[
X=|K|.
\]
Therefore
\[
L(f)=\chi(|K|).
\]

For example, if \(K\) is a filled triangle \(\Delta^2\), then \(|K|\) is contractible, so
\[
\chi(|K|)=1.
\]
Hence
\[
L(\mathrm{id})=1.
\]

\subsubsection*{Example 2: reflection of an interval}

Let \(K\) be a simplicial subdivision of the interval \([0,1]\), and define
\[
f(x)=1-x.
\]
Geometrically this is reflection about the midpoint.

The only fixed point is
\[
x=\frac12.
\]
So the fixed-point set is a single point:
\[
X=\left\{\frac12\right\}.
\]
Hence
\[
\chi(X)=1,
\]
and therefore
\[
L(f)=1.
\]

This is the simplest nontrivial example of the formula
\[
L(f)=\chi(\mathrm{Fix}(f)).
\]

\subsubsection*{Example 3: rotation of the boundary of a triangle}

Let \(K\) be the boundary of a triangle, so that \(|K|\cong S^1\). Let \(f\) rotate the triangle by \(120^\circ\).

Then no point is fixed:
\[
X=\varnothing.
\]
Therefore
\[
\chi(X)=0.
\]
Hence
\[
L(f)=0.
\]

This is a very useful example where the fixed-point set is empty.

\subsubsection*{Example 4: reflection of a square across a diagonal}

Triangulate a square by inserting one diagonal, and let \(f\) be reflection across that diagonal.

Then every point on the diagonal is fixed, and no other point is fixed. Thus
\[
X\cong [0,1].
\]
Since an interval is contractible,
\[
\chi(X)=1.
\]
Therefore
\[
L(f)=1.
\]

This is a good example where the fixed-point set is not just one point, but an entire line segment.

\subsubsection*{Example 5: antipodal map on \(S^2\)}

Consider
\[
A:S^2\to S^2,\qquad A(x)=-x.
\]
This map has no fixed points, so
\[
X=\varnothing.
\]
Thus
\[
\chi(X)=0.
\]

On the other hand,
\[
L(A)=1+\deg(A)=1+(-1)=0.
\]
So again
\[
L(A)=\chi(X).
\]

This example connects Problem 10 directly to Problem 9.

\subsubsection*{Example 6: reflection of a triangle across an altitude}

Let \(K\) be a filled equilateral triangle, and let \(f\) be reflection across an altitude through one vertex and the midpoint of the opposite side.

Then the fixed-point set is exactly that altitude segment. So
\[
X\cong [0,1].
\]
Therefore
\[
\chi(X)=1,
\]
and hence
\[
L(f)=1.
\]

This is a nice two-dimensional simplicial example in which the fixed set is visible inside the triangle.

\subsection{A compact list of the most useful examples}

The following examples are especially worth remembering:

\[
T^2\#\RP^2\cong S_3,
\]
\[
T^2\#K\cong S_4,
\]
\[
\Sigma_3\to\Sigma_2 \quad \text{as a 2-sheeted covering,}
\]
\[
S^2\to\RP^2,
\]
\[
L(\text{reflection of interval})=1,
\]
\[
L(\text{rotation of }S^1)=0,
\]
\[
L(\text{antipodal map on }S^2)=0.
\]

These examples capture the essential geometry behind Problems 7--10 in the simplest possible settings.

\end{document}
